\section{Сводка моста}
\label{sec:coords-bridge}

\begin{table}[ht]
\centering
\caption{Где живут $\symPi$, $\ln 2$ и CODATA.}
\label{tab:m-t-bridge}
\begin{tabular}{@{}p{0.42\linewidth}p{0.48\linewidth}@{}}
\toprule
Ядро $\layerMicro$ & Только мост $\layerMacro$ \\
\midrule
$\phaseVarphi_{\mathrm{disc}}$, $k_{\symPhi}$, $\symMu$; $dh$, $dt$ & $\planckHbar$, $\timePlanck$, $\lengthPlanck$, $\speedC$ \\
$\symCapitalDelta\symPhi_{\mathrm{disc}}/\phaseRingCount$ & $\symAlpha_*-1=1/(4\symPi)$ \\
$\log_2 \phaseRingCount+\varepsilon_B$ & $B_{\hv}=2\symPi/\ln 2$ \\
$\phaseSectorCount=13$ & $\lceil 4\symPi\rceil$ \\
\bottomrule
\end{tabular}
\end{table}

Спуск с макроскопической физики на~$\layerMacro$ к обратимому автомату на~$\layerMicro$ завершён.
Далее: выбор геометрии~$\carrierLattice$ (FCC), теорема дискретности и явный закон~$g$.
